\section{拓扑流形}\label{sec:topological-manifolds}

本节先研究局部形似 $\mathbb{R}^n$ 的空间, 即每个点都有一个与 $\mathbb{R}^n$ 的某个开子集同胚的邻域.

\begin{RiemannDefinition}\label{def:topological-manifold}
$n$ 维\textbf{拓扑流形} $M$ 是满足以下性质的拓扑空间:
\begin{enumerate}[label=(\roman*),leftmargin=*]
\item $M$ 是 \textbf{Hausdorff} 的, 即对 $M$ 中任意两个不同的点 $p_1,p_2$, 分别存在它们的邻域 $V_1,V_2$, 使得 $V_1\cap V_2=\varnothing$.
\item 每个点 $p\in M$ 都有一个与 $\mathbb{R}^n$ 的某个开子集 $U$ 同胚的邻域 $V$.
\item $M$ 满足\textbf{第二可数公理}, 即 $M$ 的拓扑具有可数基.
\end{enumerate}
\end{RiemannDefinition}

在定义中加入条件 (i) 和 (iii), 是为了避免这些空间的拓扑过于奇特. 特别地, Hausdorff 公理保证收敛序列的极限唯一. 这一公理与第二可数公理一起保证了单位分解的存在性 (参见第~\RGref{chap:differential-forms}~章第~\RGref{subsec:notes-integration-on-manifolds}~节). 我们将看到, 单位分解是微分几何的一项基本工具.

\begin{RiemannRemark}\label{rem:topological-manifolds-low-dimension}
若 $M$ 的维数为零, 则 $M$ 是一个赋有离散拓扑的可数集 ($M$ 的每个子集都是开集). 若 $\dim M=1$, 则 $M$ 局部同胚于开区间; 若 $\dim M=2$, 则 $M$ 局部同胚于开圆盘, 依此类推.
\end{RiemannRemark}

\begin{RiemannExample}\label{ex:manifolds-from-euclidean-space}
\begin{enumerate}[label=(\arabic*),leftmargin=*]
\item $\mathbb{R}^n$ 的每个开子集 $M$ 赋予子空间拓扑 (即 $U\subset M$ 是开集, 当且仅当 $U=M\cap V$, 其中 $V$ 是 $\mathbb{R}^n$ 的开集) 后, 都是拓扑流形.
\item (\emph{圆}) 赋有子空间拓扑的\textbf{圆}
\[
  S^1=\{(x,y)\in\mathbb{R}^2\mid x^2+y^2=1\}
\]
是 1 维拓扑流形. 条件 (i) 和 (iii) 继承自周围空间. 另外, 对每个点 $p\in S^1$, 至少有一条坐标轴不平行于 $S^1$ 在 $p$ 处的法向量 $n_p$. 投影到这条坐标轴, 就得到 $p$ 的一个 (足够小的) 邻域 $V$ 与 $\mathbb{R}$ 中一个区间之间的同胚.
\RGfigure[.46\textwidth]{circle-sphere-torus}{\textbf{a} $S^1$; \textbf{b} $S^2$; \textbf{c} 旋转环面}
\item (\emph{2 维球面}) 容易推广上一个例子, 以说明带有子空间拓扑的 \textbf{2 维球面}
\[
  S^2=\{(x,y,z)\in\mathbb{R}^3\mid x^2+y^2+z^2=1\}
\]
是 2 维拓扑流形.
\item (\emph{旋转环面}) 与前面的例子一样, 可以说明把一个圆绕一条与它不相交的轴旋转所得的旋转曲面是 2 维拓扑流形 (Fig.~\ref{fig:circle-sphere-torus}).
\item 立方体的表面是拓扑流形 (同胚于 $S^2$).
\end{enumerate}
\end{RiemannExample}

\begin{RiemannExample}\label{ex:manifolds-by-identifying-edges}
也可以用同胚把某些多边形的边等同起来, 从而得到拓扑流形. 例如, 正方形的边可以按几种不同的方式等同 (见 Fig.~\ref{fig:edge-identifications-torus-klein} 和 Fig.~\ref{fig:edge-identifications-projective-plane}):
\begin{enumerate}[label=(\arabic*),leftmargin=*]
\item (\emph{环面}) \textbf{环面} $T^2$ 是单位正方形 $Q=[0,1]^2\subset\mathbb{R}^2$ 对等价关系
\[
  (x,y)\sim(x+1,y)\sim(x,y+1),
\]
所取的商, 赋以商拓扑 (参见第~\RGref{subsec:notes-topological-manifolds}~节).
\item (\emph{Klein 瓶}) \textbf{Klein 瓶} $K^2$ 是 $Q$ 对等价关系
\[
  (x,y)\sim(x+1,y)\sim(1-x,y+1)
\]
所取的商.
\RGfigure[.52\textwidth]{edge-identifications-torus-klein}{\textbf{a} 环面 ($T^2$); \textbf{b} Klein 瓶 ($K^2$)}
\RGfigure[.52\textwidth]{edge-identifications-projective-plane}{射影平面 ($\mathbb{R}P^2$)}
\item (\emph{射影平面}) \textbf{射影平面} $\mathbb{R}P^2$ 是 $Q$ 对等价关系
\[
  (x,y)\sim(x+1,1-y)\sim(1-x,y+1)
\]
所取的商.
\end{enumerate}
\end{RiemannExample}

\begin{RiemannRemark}\label{rem:connected-sums-of-surfaces}
\begin{enumerate}[label=(\arabic*),leftmargin=*]
\item 唯一的紧致连通 1 维拓扑流形是圆 $S^1$ (见 \RGcite{Mil97}).
\item 两个拓扑流形 $M$ 和 $N$ 的\textbf{连通和}是如下得到的拓扑流形 $M\#N$: 从每个流形中去掉一个同胚于球的开集, 再用一个同胚将边界 (它们必须与球面同胚) 粘合起来 (参见 Fig.~\ref{fig:connected-sum-of-tori}). 可以证明, 任何紧致连通 2 维拓扑流形要么同胚于 $S^2$, 要么同胚于 Example~\ref{ex:manifolds-by-identifying-edges} 中的流形的连通和 (见 [\RGref{bib:Blo96}, \RGref{bib:Mun00}]).
\end{enumerate}
\end{RiemannRemark}

\RGfigure[.69\textwidth]{connected-sum-of-tori}{两个环面的连通和}
\RGfigure[.46\textwidth]{cylinder-and-mobius-band}{\textbf{a} 圆柱面; \textbf{b} Möbius 带}

若不把正方形所有的边都等同起来, 我们会得到圆柱面或 \textbf{Möbius 带} (参见 Fig.~\ref{fig:cylinder-and-mobius-band}). 这些拓扑空间是\textbf{带边界流形}的例子.

\begin{RiemannDefinition}\label{def:topological-manifold-with-boundary}
考虑\textbf{闭半空间}
\[
  \mathbb{H}^n=\{(x^1,\ldots,x^n)\in\mathbb{R}^n\mid x^n\geq 0\}.
\]
\textbf{带边界的拓扑流形}是一个具有可数开集基的 Hausdorff 空间 $M$, 使得每个点 $p\in M$ 都有一个邻域 $V$, 它或者与 $\mathbb{H}^n\setminus\partial\mathbb{H}^n$ 的某个开子集 $U$ 同胚, 或者与 $\mathbb{H}^n$ 的某个开子集 $U$ 同胚且 $p$ 对应于 $\partial\mathbb{H}^n$ 中的一个点. 第一类点称为\textbf{内点}, 其余的称为\textbf{边界点}.
\end{RiemannDefinition}

\RGfigure[.55\textwidth]{mobius-band-in-two-surfaces}{Möbius 带位于 \textbf{a} Klein 瓶内; \textbf{b} 实射影平面内}
\RGfigure[.73\textwidth]{disk-in-projective-plane}{实射影平面内的圆盘}

边界点的集合 $\partial M$ 称为 $M$ 的\textbf{边界}, 它是一个 $(n-1)$ 维流形.

\begin{RiemannRemark}\label{rem:mobius-band-decompositions}
\begin{enumerate}[label=\arabic*.,leftmargin=*]
\item 制作一个 Möbius 带的纸模型, 就容易验证其边界同胚于一个圆 (不是两个不相交的圆), 并且它只有一面 (参见 Fig.~\ref{fig:cylinder-and-mobius-band}).
\item Klein 瓶和实射影平面都含有 Möbius 带 (参见 Fig.~\ref{fig:mobius-band-in-two-surfaces}). 在射影平面上去掉这条带, 就得到一个圆盘 (参见 Fig.~\ref{fig:disk-in-projective-plane}). 换言之, 把一个 Möbius 带和一个圆盘沿各自的边界粘合, 就可以得到 $\mathbb{R}P^2$.
\end{enumerate}
\end{RiemannRemark}

若两个拓扑流形同胚, 我们就把它们视为相同. 例如, $\mathbb{R}^3$ 中半径不同的球面彼此同胚, Fig.~\ref{fig:knotted-and-unknotted-tori} 中的两个曲面也同胚. 实际上, 可以沿一个圆切开环面, 将它打结后再粘起来, 得到打结的环面. 一个显然的同胚便是将原环面上的每个点送到切开和粘合后的最终位置的映射 (不过, 这个同胚不能延拓为周围空间 $\mathbb{R}^3$ 的一个同胚).

\RGfigure[.49\textwidth]{knotted-and-unknotted-tori}{两个同胚的拓扑流形}

\begin{RiemannExercise}\label{ex:topological-manifolds}
\begin{enumerate}[label=(\arabic*),ref=\arabic*,leftmargin=*]
\item\label{ex:recognize-topological-manifolds} 下列哪些集合 (带有子空间拓扑) 是拓扑流形?
\begin{enumerate}[label=(\alph*),leftmargin=*]
\item $D^2=\{(x,y)\in\mathbb{R}^2\mid x^2+y^2<1\}$;
\item $S^2\setminus\{p\}$ ($p\in S^2$);
\item $S^2\setminus\{p,q\}$ ($p,q\in S^2$, $p\ne q$);
\item $\{(x,y,z)\in\mathbb{R}^3\mid x^2+y^2=1\}$;
\item $\{(x,y,z)\in\mathbb{R}^3\mid x^2+y^2=z^2\}$;
\end{enumerate}
\item\label{ex:classify-topological-manifolds} 上述哪些流形同胚?
\item\label{ex:klein-bottle-from-mobius-bands} 证明可以通过一个边界同胚将两个 Möbius 带粘合在一起, 从而得到 Klein 瓶 $K^2$.
\item\label{ex:connected-sum-identities} 证明:
\begin{enumerate}[label=(\alph*),leftmargin=*]
\item 对任意 2 维拓扑流形 $M$, $M\#S^2=M$;
\item $\mathbb{R}P^2\#\mathbb{R}P^2=K^2$;
\item $\mathbb{R}P^2\#T^2=\mathbb{R}P^2\#K^2$.
\end{enumerate}
\item\label{ex:euler-characteristic-of-surfaces} 2 维拓扑流形 $M$ 的一个\textbf{三角剖分}, 是将 $M$ 分解为有限个三角形 (即与 $\mathbb{R}^2$ 中三角形同胚的子集), 使任意两个三角形的交集或者是一条公共边, 或者是一个公共顶点, 或者为空 (可以证明这样的三角剖分总是存在). $M$ 的 \textbf{Euler 特征数}为
\[
  \chi(M):=V-E+F,
\]
其中 $V,E,F$ 分别为给定三角剖分的顶点数, 边数和面数 (可以证明它是良定义的, 即不依赖于三角剖分的选择). 证明:
\begin{enumerate}[label=(\alph*),leftmargin=*]
\item 给三角剖分增加一个顶点不会改变 $\chi(M)$;
\item $\chi\bigl(S^2\bigr)=2$;
\item $\chi\bigl(T^2\bigr)=0$;
\item $\chi\bigl(K^2\bigr)=0$;
\item $\chi\bigl(\mathbb{R}P^2\bigr)=1$;
\item $\chi(M\#N)=\chi(M)+\chi(N)-2$.
\end{enumerate}
\end{enumerate}
\end{RiemannExercise}
